\section{An exact obstruction at bad coefficient primes}
\label{sec:bad}

The full \emph{integral} matrix-span hypothesis is essential. Absolute
irreducibility over the fraction field is too weak, even in rank two.

\begin{theorem}[Extra torsion invisible to the trace depth]
\label{thm:bad}
Fix an odd prime \(p\) and positive integers \(a,c\). Let \(O=\mathbb Z_p\)
and \(G\) be free on \(u,v,z\). Define
\[
 \begin{aligned}
 \rho(u)=\sigma(u)&=I+p^cE_{12},&
 \rho(v)=\sigma(v)&=I+p^cE_{21},\\
 \rho(z)&=I,& \sigma(z)&=(1+p^a)I.
 \end{aligned}
\]
Both representations are absolutely irreducible over \(\mathbb Q_p\).
Their trace-congruence depth is exactly \(a\), but
\begin{align}
 \operatorname{length}_O
 \operatorname{Hom}_{(O/p^n)[G]}(\rho_n,\sigma_n)
 &=\min(n,a)+3\min(n,a,c), \label{eq:badlength}\\
 \sup_n\operatorname{length}_O
 \operatorname{Hom}_{(O/p^n)[G]}(\rho_n,\sigma_n)
 &=a+3\min(a,c).
 \end{align}
The image order generated by \(u,v\) is
\[
 \mathcal A=OI+Op^cE_{12}+Op^cE_{21}+Op^{2c}E_{11}.
 \]
It has colength \(4c\) in \(M_2(O)\) and conductor exponent \(2c\):
\(p^{2c}M_2(O)\subseteq\mathcal A\), but
\(p^{2c-1}M_2(O)\nsubseteq\mathcal A\).
\end{theorem}

\begin{proof}
Over the algebraic closure the first unipotent has unique invariant line
spanned by the first coordinate vector, and the second has unique
invariant line spanned by the second. They have no common line.

The representations agree modulo \(p^a\) on generators and hence on
all group elements. At \(z\) the trace discrepancy is \(-2p^a\),
of valuation \(a\).

Write an intertwiner as
\[
 X=\begin{pmatrix}t&y\\z'&t+h\end{pmatrix}.
 \]
Commutation with the first two matrices imposes exactly
\(p^cy=p^cz'=p^ch=0\) modulo \(p^n\).
The third equation imposes \(p^aX=0\), killing all four coordinates.
Their lengths are \(\min(n,a)\) and three copies of \(\min(n,a,c)\).

The displayed order is closed under products and contains the generators
and their inverses. Conversely their differences from \(I\) give the
off-diagonal generators, and their product gives \(p^{2c}E_{11}\).
Its elementary divisor valuations are \(0,c,c,2c\), proving the
colength and conductor statements.
\end{proof}

For \(a\geq c\), the excess over trace depth is exactly \(3c\).
For instance \(p=3,a=2,c=1\) gives trace depth \(2\) and global
intertwiner length \(5\). This is not a sampling problem: even the
traces on \emph{all} group elements fail to determine the extra torsion
by the full-span formula. An integral model or an explicitly measured
defect of its image order must accompany any extension to bad primes.

The relation matrix has Smith factors
\[
 p^{\min(a,c)},\quad p^{\min(a,c)},\quad
 p^{\min(a,c)},\quad p^a .
\]
The checker verifies them for four parameter choices and independently
enumerates all solutions for small quotients. The general formula
follows from the proof, not from extrapolating the computations.
